\documentclass[11pt]{amsart}
\usepackage[T1]{fontenc}
\usepackage[utf8]{inputenc}
\usepackage{amsmath,amssymb,amsthm,mathtools}
\usepackage{microtype}
\usepackage[colorlinks=true,linkcolor=blue,citecolor=blue,urlcolor=blue]{hyperref}

\newtheorem{theorem}{Theorem}
\newtheorem{lemma}[theorem]{Lemma}

\theoremstyle{remark}
\newtheorem{remark}[theorem]{Remark}

\newcommand{\Cyc}[1]{C_{#1}}
\newcommand{\Z}{\mathbb{Z}}
\newcommand{\ord}{\operatorname{ord}}


\begin{document}

\title[No Cyclic $(9423,4711,2355)$ Difference Set]{On the Nonexistence of a Cyclic $(9423,4711,2355)$-Difference Set\thanks{AutoMath assisted in generating the solution and preparing this draft.}}
\date{April 15, 2026}

\begin{abstract}
We prove nonexistence of a cyclic difference set with parameters $(9423,4711,2355)$.
The proof uses the multiplier $2$ and the quotient of order $3$.
After translating a hypothetical difference set so that it is fixed by doubling, the two nonzero quotient classes mod $3$ have equal size.
The corresponding order-$3$ character value is therefore an integer, but its square must equal $2356$, which is impossible.
\end{abstract}

\maketitle

\section{Introduction}
Baumert and Gordon list $(9423,4711,2355)$ among the small open cyclic Hadamard cases \cite{BG2004}, and Gordon's La Jolla status tables continue to retain $9423$ on the residual list \cite{Gordon2019}.
We record a one-page nonexistence proof.

\begin{theorem}
The cyclic group $\Cyc{9423}$ admits no $(9423,4711,2355)$-difference set.
\end{theorem}

\begin{proof}
Assume for contradiction that $D\subseteq G=\Cyc{9423}$ is a $(9423,4711,2355)$-difference set.
Then
\[
n=k-\lambda=4711-2355=2356.
\]
Since $2$ divides $n$ and $\gcd(2,9423)=1$, the multiplier theorem implies that $2$ is a multiplier.
Thus there exists $t\in G$ such that
\[
2D=D+t.
\]
Because
\[
\gcd(2-1,9423)=1,
\]
we may translate $D$ and assume from the start that
\[
2D=D.
\]

Let $H=3G$, so $|H|=3141$ and $G/H\cong \Cyc{3}$.
Write
\[
x_0=|D\cap H|,\qquad
x_1=|D\cap (1+H)|,\qquad
x_2=|D\cap (2+H)|.
\]
The doubling map induces multiplication by $2$ on $G/H\cong \Cyc{3}$, so it swaps the two nonzero classes.
Since $2D=D$, we obtain
\[
x_1=x_2.
\]

Let $\psi$ be a nontrivial character of $G/H$, and write $\omega=\psi(1+H)$, so $\omega^2+\omega+1=0$.
Lift $\psi$ to a nonprincipal character $\chi$ of $G$.
Then the standard character identity for difference sets gives
\[
|\chi(D)|^2=n=2356.
\]
On the other hand,
\[
\chi(D)=x_0+x_1\omega+x_2\omega^2 = x_0+x_1(\omega+\omega^2)=x_0-x_1,
\]
because $x_1=x_2$ and $\omega+\omega^2=-1$.
Thus $\chi(D)$ is an integer.
Consequently $2356=|\chi(D)|^2$ would be the square of an integer.
But
\[
2356 = 4\cdot 19\cdot 31
\]
is not a perfect square.
This contradiction proves the theorem.
\end{proof}

\begin{remark}
The obstruction is already visible on the index-$3$ quotient.
Once the $2$-multiplier normalization is in place, the order-$3$ character sum becomes integral, while the forced norm $2356$ is nonsquare.
\end{remark}

\begin{thebibliography}{99}

\bibitem{BG2004}
Leonard D.~Baumert and Daniel M.~Gordon,
\emph{On the existence of cyclic difference sets with small parameters},
Fields Inst. Commun. \textbf{41} (2004), 29--40.

\bibitem{Gordon2019}
Daniel M.~Gordon,
\emph{The La Jolla Difference Set Repository},
ArasuFest talk slides, August 3, 2019.

\end{thebibliography}

\end{document}
